


\documentclass[11pt, twocolumn]{article}


\usepackage[a4paper,
            bindingoffset=0.2in,
            left=.35in,
            right=.35in,
            top=.5in,
            bottom=.5in,
            footskip=.35in]{geometry}

\setlength{\columnsep}{38pt}

%\titleformat{\section}
%  {\normalfont\Large}{\thesection}{1em}{}

\renewcommand{\thesection}{\textcolor[rgb]{0,0.7,0.4}{\arabic{section}}}%


%\titleformat{\subsection}
%  {\normalfont\large}{\thesubsection}{0}{}

%\definecolor{subsclr}{rgb}{0,0.5,0.7}
%\definecolor{es-aqua}{HTML}{00ffffff}% 

\renewcommand{\thesubsection}{\textcolor[rgb]{0,0.2,0.4}{\arabic{section}}%
\hspace{1.5pt}.\hspace{1.5pt}\small{\textcolor[rgb]{0,0.3,0.4}{\arabic{subsection}}%
}}

\usepackage{xparse}

\makeatletter

%\newcommand{\@sect}[1]{\section{#1}}
%\newcommand{\@subsect}[1]{\vspace{9pt}\subsection{#1}}

\newcommand{\sL@section}[1]{\section{#1}}
\newcommand{\sL@subsection}[1]{\vspace{14pt}\subsection{#1}

\vspace{-18pt}
\textcolor[rgb]{0.6,0.6,0.8}{\rule{\columnwidth}{0.4pt}\vspace{-15pt}}}

\newcommand{\@sectionLevel}[2]{%
\ifcase#1\relax
\or{\sL@section{#2}}
\or{\sL@subsection{#2}}
\else{}
\fi
}


\NewDocumentCommand{\s}{E{.}{}D||{}m}{%
\IfNoValueTF{#1}{
\@sectionLevel{#2}{#3}
}
{%else
\@sectionLevel{#1}{#3}
}
}


\newcommand{\@pLevel}[2]{%
\ifcase#1\relax{p0=#2}
\or{\vspace{6pt}#2}
\or{p2=#2}
\else{}
\fi
}


\NewDocumentCommand{\p}{E{.}{}D||{}+m}{%
\IfNoValueTF{#1}{
\@pLevel{#2}{#3}
}
{%else
\@pLevel{#1}{#3}
}
}


\newcommand{\@exsRef}[1]{#1}

%\usepackage{suffix}
%\WithSuffix\newcommand{\exsRef*}{:#1}

\NewDocumentCommand{\exsRef}{sd()}{%
\IfBooleanTF{#1}{\@exsRef#2}{(\@exsRef#2)}}

\makeatother


\newcommand{\citeLabel}[1]{\cite{#1}}

\NewDocumentCommand{\citePage}{md()}{%
\cite[p. #2]{#1}}

\NewDocumentCommand{\citePages}{md()}{%
\cite[pp. #2]{#1}}

\NewDocumentCommand{\citeLocator}{md()}{%
\cite[#2]{#1}}


\newcommand{\onethrft}{(1)--(14)}


\newcommand{\eS}[1]{\textbf{\scriptsize{#1}}}
\newcommand{\q}[1]{``#1''}

\newcommand{\mdash}{---}

\newcommand{\visbreak}{\vspace{9pt}}




\newcommand{\upAdj}[1]{\vspace{-#1}}


\newcommand{\nip}{\noindent}

\newcounter{exsCount}

\newcommand{\exsItem}{%
\stepcounter{exsCount}
\noindent(\theexsCount) 
}

\newenvironment{exsGroup}{\vspace{5pt}}{\vspace{5pt}}

\newenvironment{enums}{\vspace{5pt}\begin{enumerate}}{\end{enumerate}}

\newenvironment{itemz}{\vspace{5pt}\begin{itemize}}{\end{itemize}}


\usepackage{csquotes}
\newenvironment{blockQuote}{\vspace{5pt}\begin{displayquote}}{\end{displayquote}}

\newcommand{\enumsItem}{\item}
\newcommand{\symItem}{\item}

\usepackage{textcomp}
\newcommand{\sarr}{\textrightarrow}

\newcommand{\OxFO}{\textbf{\scriptsize{0xF0}}}

\newcommand{\UTFonesix}{\textbf{\scriptsize{UTF-16}}}

\newcommand{\etal}{\textit{et al.}}

\newcommand{\Gardenfors}{G\"ardenfors}



\newcommand{\redbox}{\makebox{\eS{red}\hspace{4pt}\ensuremath{\cap}\hspace{4pt}\eS{box}}}

\newcommand{\jrefs}{(15)/(16)}

\newcommand{\cq}[1]{{\fontfamily{gar}\selectfont ``}#1{\fontfamily{gar}\selectfont ''}}

\newcommand{\repeatName}{\_\_\_\_\_}

\newcommand{\visavis}{vis-`a-vis}


\usepackage{enumitem}
\setlist[enumerate]{leftmargin=4pt,labelindent=-3pt}
\setlist[itemize]{leftmargin=4pt,labelindent=-3pt}


\usepackage{xcolor}



\usepackage[tracking=true]{microtype}

\newcommand{\tinyurl}[1]{\raisebox{2pt}{{\scriptsize \url{#1}}}}

\usepackage[colorlinks=true]{hyperref}

\colorlet{urlclr}{red!40!magenta!50!orange}

\hypersetup{
 urlcolor = urlclr,
 urlbordercolor = cyan!60!black,
 linkcolor = red!30!black,
 citecolor = orange!30!black,
 citebordercolor = yellow!30!black,
} 

\usepackage{setspace}

\newcommand{\biburl}[1]{ {\fontfamily{gar}\selectfont{\textcolor[rgb]{.2,.6,0}%
			{\scriptsize \\\setstretch{3}{\textls*[-70]{\url{#1}}}}}}}

%\newcommand{\bibformat}[1]{ {\fontfamily{gar}\selectfont{\textcolor[rgb]{.6,.2,0}%
%			{\scriptsize \\\setstretch{4}{{#1}}}}}}

\newcommand{\bibformat}[1]{ {\fontfamily{gar}\fontsize{10}{10}\selectfont{%
\textcolor[rgb]{.6,.2,0}{%
\textls*[-70]{%
#1%
}}}}}



%\usepackage{verbatim}
%\usepackage{verbatimbox}

\usepackage{listings}
\lstset{
basicstyle=\ttfamily\linespread{0.94}\fontsize{8}{11}\selectfont\color[rgb]{.6,.2,0},
columns=flexible,
breaklines=true
}

\newsavebox\hrefbox

\NewDocumentEnvironment{boxurl}{d()}{%

\begin{lrbox}{\hrefbox}
\begin{minipage}{\columnwidth}}
{\end{minipage}
\end{lrbox}

\vspace*{-4pt}
\href{#1}{\usebox{\hrefbox}}}


\usepackage{titlesec}


\setcitestyle{square}

\usepackage[T1]{fontenc}
\usepackage{tgtermes}

\raggedbottom

%\newgeometry{left=0.8in,right=0.8in,top=.6in,bottom=.6in}



\usepackage{natbib}
\setlength{\bibsep}{3pt}








\begin{document}


\begin{abstract}
This paper examines methodological foundations for
linguistics in terms of phenomenology and the
philosophy of science.  I argue for explanatory
paradigms that adequately model the interdependence of
linguistic and extra-linguistic knowledge/faculties
for understanding typical sentences.  In lieu
of set theory or preducate logic formalization, I
point out some similarities between natural and computer
languages, so that formal structures related to
compilers and programming-language implementation
may offer useful simplified models for human
communication.  In particular, I emphasize how
surface discourse may be transformed to sequences
of operations toward conceptual aggregation/synthesis,
while maintaining that most details of how
concepts are pairwise combined, modifying one another,
depend on extra-linguistic cognitive activity.
Situational meaning construal is outside the
scope of linguistics by analogy to how
system kernels and Central Processing Units
are outside the scope of programming language theory.
\end{abstract} 


\s|1|{Introduction}

\p.1{%
It is clearly true that an important part of human language is governed by rules or
structures, ones that can be (at least to a useful extent) formalized and computationally
simulated.  And yet, language was not designed by committee: claiming that a
particular semantic or syntactic pattern has become entrenched is not a
complete argument, as if language were intrinsically a collection of formal rules, and we
need merely infer them (like reverse-engineering compiled computer code).  The
question is \textit{why} patterns become entrenched and stay so.
} % end paragraph 

\p.1{%
Plus, we should not
overestimate rules' rigidity.  That a bachelor is an unmarried man, for example,
is a canonical example of logic underlying semantics; but sentences like
these are quite meaningful:

\begin{exsGroup}

\exsItem{}   My brother has been married for 10 years, but in many ways he's still a bachelor.

\exsItem{}   All the eligible bachelors in this town are married.

\end{exsGroup}

\nip
Moral of the story: just as, in a fair legal system, it should be assumed that
almost all laws have instances where reasonable people would break them, so
in language, for every \q{rule} there are contexts where competent
speakers will ignore them.
} % end paragraph 

\p.1{%
To motivate further discussion about the limits of linguistic
formalization, consider these pairs of sentences:

\begin{exsGroup}

\exsItem{}   Everyone sang along to three songs

\exsItem{}   All appeals are heard by a panel of three judges

 \visbreak{}

\exsItem{}   All New Yorkers live in one of five boroughs

\exsItem{}   All New Yorkers complain about how long it takes to commute to New York City.

 \visbreak{}

\exsItem{}   In Ariel's picture of Beatrice, she looks like she has a cold

\exsItem{}   In Ariel's picture of Beatrice, she snapped the lens just as two cardinals flew by

 \visbreak{}

\exsItem{}   We operated near the bank

\exsItem{}   On the mantle was a wooden bat

\end{exsGroup}

\nip
The first three pairs have interpretive details apparently reliant on extra-linguistic
reasoning.  The situation in \exsRef(3) and \exsRef(4) is quantifier scope: assuming \exsRef(3) discusses some
music performance, it would be abnormal for everyone in the audience to distinctly choose
exactly three songs to sing along with.  Hence, we can assume that, among all songs
played that evening, there were specifically three with the property that (all) the audience
joined in.  The numeral-quantifier \q{three} is thus outside the scope of \q{everyone}.
By contrast, one is free to hear \exsRef(4) as nesting the scope of \q{three} inside \q{All}
(it is a general principle that three judges hear an appeal, but the set of
individuals can vary from one appeal to another).
} % end paragraph 

\p.1{%
In \exsRef(6), interpreting a quanitifier's intent requires distinguishing levels of exactitude.  Rhetoriccally,
\q{all} often actually means \q{most}.  The \exsRef(6) speaker is not heard to
report on the attitude of 10+ million residents as if they were personally interviewed (nor would \exsRef(6) be
falsified by scattered counter-examples).  On the other hand, the land area of New York is indeed
precisely divided into five boroughs, so \exsRef(5) is \textit{literally} correct \mdash{} but only via a rigid
reading of the referential intent behind \q{New Yorker}.  In many contexts this term extends
to many beyond the city proper.
} % end paragraph 

\p.1{%
Uncertainties in \exsRef(7) and \exsRef(8) are of a different sort: anaphora resolution rather than quantifier-scope.
They also have two competing interpretations, depending on whether \q{she} is read as Ariel or Beatrice.  The
former option is more likely in \exsRef(7) because one can tell from a photograph what the depicted subject
"looks like".  But we can't \textit{a priori} exclude a continuation like \q{you'd have to be sick to
take such a bad shot}.  Nor can we fail to imagine a scenario for \exsRef(8) where Ariel takes a picture of
Beatrice photographing colorful birds.
} % end paragraph 

\p.1{%
The issue in \exsRef(9) and \exsRef(10), by contrast, is ordinary polysemy: \exsRef(9) might involve an
accountant with offices near a financial institution, or a field clinic established
near a river.  And \exsRef(10) could describe a baseball collectible or an animal sculpture.
I include these for point of contrast: the ambiguity in the former three pairs,
I believe, is \textit{holistic} and \textit{compositional} whereas the latters' underdeterminism
is more localized (even though selecting one or another reading of \textit{bank}, and
\textit{bat}, frames lexical choices for the overall sentense; my point is that
the locus of sense-divergence can be reduced to a single word).  Hold that thought.
} % end paragraph 

\p.1{%
I chose these sentence-pairs because the variable factor in each case flip-flops from first to second, at least
in the most \textit{likely} reading.  But subtle variation can point to the opposite, as in these alternatives
for \exsRef(3) and \exsRef(4):

\begin{exsGroup}

\exsItem{}   Every band performed three songs

\exsItem{}   Everyone in the audience sang along to one song or another

\end{exsGroup}

\nip
Changing the context to a \textit{band} playing, versus an audience member, imposes a different interpretive
framework: usually bands perform a \q{set} with their own material (or covers).  Meanwhile, switching
from the exact \q{three} to the more flexible \q{one or another} gives enough construal slack that
putting the later scope in the former makes interpretive sense.
} % end paragraph 

\p.1{%
Despite the competing interpretations, each reading of each sentence I have discussed so
far, I believe, has a fairly obvious logical form.  But there are other examples where a logical
gloss would not fully capture the communicated content of a given sentence.  Consider:

\begin{exsGroup}

\exsItem{}   This election is slipping away from the Democrats

\exsItem{}   Student after student came by to complain about the reading assignments

\end{exsGroup}

\nip
Notice how \exsRef(13) and \exsRef(14) subtlely describe warrants as well as assert beliefs.
The metaphor of \q{slipping away} connotes a
gradual process that the speaker has observed from an extended period of time.  Similarly,
\exsRef(14) suggests multiple similar episodes that collectively give a warrant,
presented in language that suggests an extended,
enduring process, expressing a kind of \q{narrative account} of the speaker's evidential
rationale.  The details are different than if all objections were made during one
heated class session, say.
} % end paragraph 

\p.1{%
With no further qualification, of course, we
hear an implicit \q{I believe that} attached to any utterance.  Since a
claim of \eS{P} is almost always a claim of \q{I believe \eS{P}} this kind of
meta-discursive attitude can be subsumed into ordinary logic.
Often, of course, we present one claim to justify another, as in
(a variant form of \exsRef(11), say, might explicitly mention polling numbers).
So, in short, the pattern of propositions justifying other propositions
does not in itself challenge the validity of semantic theories
premised on treating sentences as encodings for propositions.
} % end paragraph 

\p.1{%
However, the resources a speaker has for expressing the backstory of their
propositional attitudes span a wide spectrum of connotations and experiences,
and I'd argue that there is no logical system that could precisely model
how language-users present details related to how one belief compels
or reinforces another.  Most discourse embeds propositions not in
chains of deductive reasoning but rather in narrative presentations
that reflect the passing of time, the interlocking of situations, and
in general the somewhat informal manner in which
we formulate beliefs and opinions.  To the degree that we
express beliefs via narrativizing constructions or metaphors like
\q{slipping away} and \q{student after student} \mdash{} or so I claim \mdash{} this
part of the discourse cannot be adequately captured via logic, but
rather depends on lifeworld processes responsible for forming opinions.
} % end paragraph 

\p.1{%
Cumulatively, then, I believe that sentences \onethrft{} offer examples
of two issues for formal semantics:

\begin{enums}

\enumsItem{}   While assertion (or implicature) of some proposition \eS{P} is a
feature of most sentences, there may be details contained in
the meaning of the sentence that are not contained in \eS{P} itself
(or in logically related propositions, such as \textit{belief that \eS{P}}). 
\enumsItem{}   Even if the meaning of sentence \eS{S} is indeed to a large extent included in some
\eS{P,} background knowledge may be needed to infer \eS{P} as the intended
meaning.  There may be no deterministic process to extract \eS{P} from
\eS{S} via syntactic, semantic, or pragmatic rules alone.

\end{enums}

\nip
If these points are valid, then the project of formulating semantics as the search
for rules governing how sentences \eS{S} encode propositions \eS{P} is notably limited.
} % end paragraph 

\p.1{%
I don't want to make too much out of ambiguity in and of itself.
A \q{dresser} can be someone donning clothes or a piece of furniture.
So people must discern which word-sense applies, usually with the aid of background knowledge.
Anyone who believes that such facts problematize a particular semantic theory
need not progress to \textit{sentence-level} ambiguity.
} % end paragraph 

\p.1{%
However, the \textit{lexical} underdeterminacy of \q{dresser} or \q{bat}
is a purely local phenomenon that can't be equated with more holistic
examples such as \onethrft{}.
} % end paragraph 

\p.1{%
To put it differently, it's credible to imagine homonyms like \q{dresser}
as surface-level accidents, as if we can avail of a
lexical \q{deep structure} by analogy to the syntactic.
Perhaps lurking behind English lies an idealized language where senses
could be marked, so \q{person} and \q{furniture} dresser become just
different words.  The problem on a larger scale is that sentences are supposed to
function via composition: as words signify concepts, so sentences
convey propositional content because they merge concepts according
to specific recipes (syntax matters because multi-concept complexes
are not just unstructured \q{bags}: the King loves the Queen versus vice-versa).
I argue from cases like \onethrft{} that even after fully digesting
compositional intent, predicate signification remains underdetermined.
} % end paragraph 

\p.1{%
To borrow the apologia from the lexical realm, we'd have to argue that
lurking behind any actual dialect is an idealized language which
\q{marks} all logically salient elements \mdash{} quantifiers, anaphora, deictics
(imagine marking a proximate versus distal \q{there}, or abstract versus concrete
\q{here}),
conjunctions (cf. \q{and} as temporal as in \q{John quit his job and moved out of town},
wherein we hear \q{and \textit{then}}) \mdash{} with enough granularity to forestall
all compositional ambiguities.  We can imagine some inflectional variation
on words like \q{all} to distinguish crisp from fuzzy domains, or nested from
independent quantification; or morphological features on pronouns to
distinguish verb subjects from objects in anaphoric reference.  But
should we theorize ideal languages that subsume all such ambiguities
under fully-determined contracts?
} % end paragraph 

\p.1{%
We might be tempted by a formal semantcs whose basic premises are like these:

\begin{enums}

\enumsItem{}   Canonically, sentences signify propositions \mdash{} by analogy to how words signify concepts.
 
\enumsItem{}   Behind any actual language is an idealized language wherein sentences denote
structured propositions via compositional patterns, themselves encoded through
syntactic conventions and the semantic properties of specific word-classes
(quantifiers, pronouns, and so forth).

\end{enums}

\nip
To the degree that speakers need to employ extralinguistic reasoning,
those occasions would then be treated as individuals mentally refining the
actual lexicon to an approximated ideal.
It is true, say, that \q{Trump} should be interpreted differently
when people are playing bridge or talking politics, but this means that
ambient situations nudge speakers in contextual ways when construing how
their spoken language should be projected to its perfected form.  Context
disposes us to \q{label} a word like \q{Trump} as either a person or a card
suit.  But such are antecedent to semantics proper; for theoretical
purposes we can take the lexicon as enriched with such labels \textit{a priori}.
The cognitive gap between actual and ideal dialects is orthogonal to linguists'
project, which is to describe rules of the latter.
} % end paragraph 

\p.1{%
I contend that paradigms authorizing linguists and philosophers to
imagine such ideal languages \mdash{} as their focal topic \mdash{} leads us
to underestimate the theoretical dependence of linguistics
on other disciplines.  By way of example, sentences 1-12 suggest
that the specific \eS{P} derived from an \eS{S}
depends (in part) on addressee's rational acts that aren't linguistic
in nature.  The reason why quantifier scope is nested in \exsRef(4) but not \exsRef(3) is that
the predicates signified by that reading make more sense than their
alternatives, but such judgment relies on knowledge of socio-cultural, not
linguistic, facts.  In short, whatever a sentence means, it does so
\mdash{} in the general case \mdash{} not via linguistic structures alone, but
via the interplay of such structures with more general cognitive
and social-semiotic processes.  Consequently, linguistic
explanation alone is incomplete.  Moreover, these issues are sufficiently
holistic and entrenched that we cannot dismiss them as cognition
just bridging a gap between real and ideal language, by analogy to
real and ideal dictionaries (where each distinct word-sense might
be marked in some disambiguating manner).
} % end paragraph 

\p.1{%
No prehistoric King ever instructed his subjects to go create a
language for expressing truth conditions; even when \q{the meaning of a
proposition is given by its truth
conditions} we have to ask \textit{why} sentences encode truth conditions (or
are optimized to do so).  People don't talk out of some abstract desire
to assert fragments of symbolic logic, but rather to coordinate their joint
activity.  It's easy to see why logically transparent statements are an
important part of that, but still, language's ability to convey beliefs
is an \textit{emergent property} of human verbal communication.
The existence of coherent \eS{S-to-@}/P mappings is not an axiom to
deploy in linguistic explanations, but a cognitive faculty to be explained.
} % end paragraph 

\p.1{%
To be clear, I do believe the value of signifying propositions is
self-evident.  My point is that any instance of \eS{S} conveying \eS{P}
is not just a linguistic phenomenon, but rather a detail of
two (or more) persons' social and collaborative interactions.
If I slip, fall, clutch my ankle, and grimace, someone will
conclude that I might be injured; but groaning \q{ouch} does not
express this through linguistic means.  Similarly, if I see a
car veering dangerously close to my friend John from behind him, shouting
\q{John} to get his attention does not make that word signify
\q{watch out!}.  In most communication, linguistic structure
plays a more prominent role \mdash{} only rarely do we convey ideas through vocalization
alone, untethered from syntax and semantics (and pragmatics) \mdash{}
yet just because propositions get conveyed
in a linguistic context does not ensure that \textit{why} or \textit{how} this
happens is fully explainable in linguistic terms.
} % end paragraph 

\p.1{%
Lexical plurality is one domain where
linguistics is obviously incomplete, because background knowledge
is needed to adjudicate competing word-senses.  But I believe
the \textit{compositional} indeterminacy of whole sentences
is qualitatively different, and reveals a more substantial
surface area where linguistic and extralinguistic
phenomena inter-penetrate.  Hence, linguistics
must actively connect its paradigms and presuppositions
to a larger intellectual environment, where
overarching theories of cognition and communication
are formulated and disputed.
} % end paragraph 

\p.1{%
For these reasons, I believe \textit{cognitive phenomenology} is
relevant to linguistics because it potentially provides this
overarching theoretical framework.  Here, I will outline
why I think Husserlian phenomenology can play this role,
and then what implications for syntactic theory might follow
therefrom.
} % end paragraph 


\s|1|{Phenomenology and the Philosophy of Science}

\p.1{%
I argue in this section that (Husserlian) phenomenology and the philosophy of
science can jointly offer a framework for justifying linguistic
methodology.  On (overly superficial, I believe) interpretations,
phenomenology is often seen as hostile toward science.  I propose, however,
that phenomenology's analytic role relative to other humanistic
disciplines can be usefully discussed in terms of the mutual support
between interrelated scientific theories.
} % end paragraph 

\p.1{%
An intrinsic aspect of any science is establishing a collection of
entities (or kinds thereof) that comprise its focus.  Apart from
elementary particle physics (perhaps) such entities will almost
always be reducible to smaller parts.  Hydrology depends on
properties of water molecules, and botany on plant cells.
However, the discourse of the higher-scale disciplines should
not parrot that of chemistry or cellular biology.  It is the
responsibility of each higher-scale theory to develop a
vocabulary and classification that describes dispositions
and attributes of their higher-scale objects.  Explanations can
then address why posited dispositions and attributes
generate observed behaviors.
} % end paragraph 

\p.1{%
By analogy, there is a distinct set of terms relevant to
high-level computer source code: procedures, lexical variables,
lifetimes, types, pointers, execution branches.  Separately,
there are terms related to machine code: registers, instruction
sets, opcodes, memory pages.  Although source files compile
to machine code, analyses of programming languages' rules and
implementations cannot in general \q{quantify} over machine-level
concepts.  Moreover, it would be impossible to explain the rational
structure behind machine-code sequences without employing terms
or concepts derived from the original high-level sources.  Examining progressively smaller
physical scales may suffice for purely mechanical explanations of
physical processes, but remains explanatorily incomplete insofar
as macro-scale phenomena structure and mold the micro-scale aggregates.
By analogy, a purely molecular account of the dynamics of a spinning
baseball does not take into consideration the thought processes
of the pitcher who decided to throw a curveball.  A complete
account of the universe in that moment needs to recognize such
work of mind as part of the causal background for the ball's spin.
} % end paragraph 

\p.1{%
In general, causative explanation is incomplete on any one scale in
particular; it is only by bridging multiple theories that one can
simultaneously analyze both purely physical properties (dependent
on micro-scale realities) and aggregate, complex behavior
(caused by interactions between higher-scale objects).  Individual
theories contribute to cross-scale explanations by formalizing
their own terminological and epistemological frameworks, ones
that focus on some class of entities but also explicate how
peer theories (e.g., those of smaller constituents) describe
causal grounds for the dispositions and attributes studied by
the original theory.  This is intrinsic to any science: hydrology
and botany (reusing those examples) exist as such because it
is possible to negotiate discursive practices and investigatory paradigms
for macro-scale complexes of water (rivers, oceans, lakes, dams)
and \textit{Plantae} organisms (trees, forests, foliage ecosystems, etc.) respectively, all the
while connecting such paradigms to theories of smaller constituents
(water molecules and plant cells).
} % end paragraph 

\p.1{%
The situation is similar with respect to humanities and social sciences.
A political theorist might note the waxing and waning of parties' popularity,
but such phenomena ultimately rest on individual voters' decisions.
The mindset of a single person acts as a societal \q{reductive base}
(likewise for economics and the individual worker/consumer).  Such brings
one naturally to cognition, which is an inevitable explanatory
setting for any talk of decision-making or ideology.
} % end paragraph 

\p.1{%
Inter-theory reduction depends on filling macro-to-microscale gaps with mid-level
analyses; source files, for example, become transformed to several
intermediate representations during a compiler pipeline.  Vis-a-vis the
human mind, neurologists certainly understand how clusters of neurons
form \q{nuclei} or \q{ganglia} with distinct functional roles; and,
on a higher scale, brain scientists identify many distinct
(and functionally unique) anatomic regions.  Working from
\q{the ground up}, however, one is still far from connecting
neurophysical substrata to observed patterns in human cognition.
} % end paragraph 

\p.1{%
There is an almost universally recognized explanatory gap with respect
to consciousness itself: we still have little understanding of
how purely biological phenomena can give rise to subjective
\q{first-personal} states and awareness.  Some philosophers
apparently minimize this gap by construing consciousness as epiphenomenal,
ontologically separated from the primary functional structures
of cognition.  However, the nature of consciousness \textit{per se} clearly
plays a critical role in the mental careers of conscious animals.
Consciousness evolved at some point in natural history, and presumably
it conferred adaptative benefits.  Instead of reacting to stimuli
merely via instinct (apart from rare
\q{fight or flight} reflexes), forming conscious mental representations
\mdash{} with a certain distance between observer and
observed \mdash{}
evidently permits a wider and more malleable repertoire
of responses.  Animal behavior and intelligence
has evolved in consort with consciousness: we exist in a world
of conscious representations that allow us to observe
surrounding environments without immediately acting.  We benefit
from the opportunity to formulate reactions thoughtfully, informed
by values, personal identity, and past experience.
} % end paragraph 

\p.1{%
Suppose I stand by a platform waiting for a subway to arrive.  I need to board
the correct train.  To do so, I may well follow a well-rehearsed
protocol: hearing the train cues me to look down the platform; upon seeing
its outline I would then focus attention on the letter- and color-coded
symbol painted on the front.  If it indicates the correct train,
I stepp back to allow passengers
to disembark.  Each step in this process is consciously mediated.
I need to deliberately shift my perceptual orientation multiple
times, and could break off the \q{protocol} given some external
stimuli.  Perhaps I hear my cell phone pinging.  Or, there's an announcement
that the current train is out of service, so it's a moot point.  The key
detail is that such protocols don't just occur by instinct.  They continue,
or not, because of subjectively controllable choices of cognitive orientation.
} % end paragraph 

\p.1{%
Not that we have first-person awareness of all cognitive
processes.  Conscious reality has many constituents
that we experience as perceptual givens despute being the product of
complex (subconscious) reasoning.  We see a world of discrete objects,
whose contours and boundaries are inferred via visual processing.
Sight employs color variation to partition visual scenes, but
it is not as simple as isolating distinct color patches.  Given a
red rose in a blue vase, we cannot segment the image merely by
treating red as the flower and blue as the pottery.
Color-perception has added complexity first because some chromatic
variation is expected in an extended region due to texture
and lighting; and second because visible substances have
optical qualities beyond just red/green/blue quantities
\mdash{} albedo, specularity, glossiness, metallicity, opaqueness, etc.
Object boundaries are suggested by discontinuities in visual properties of
varying kinds, not just pure color; moreover, segmentation
algorithms have to be cautious of spurious boundaries
due to light/shadow effects and the vagaries of pigmentation.
} % end paragraph 

\p.1{%
Humans subconsciously perform complex image-analysis tasks, many
of which can be simulated by computer-vision algorithms:
texture segmentation, 3D reconstruction, occlusion compensation
(e.g., perceiving the shape of a vase's back side by projecting
around from its visible rim-circle).  Via color alone we
infer textural and optical details, yielding a significantly
enlarged color space, and thereby perceive objects' external
geometry (e.g., their gradient vectors at each point
\mdash{} the degree to which their shape departs from a flat plane).
We do not consciously \textit{direct} these processes, by contrast to
my example of ascertaining which train is arriving.  Nevertheless,
conscious experience is intrinsically structured by
faculties of visual (and, to some extent, haptic/tactile)
perception so that we can trust our senses as a usually
reliable source of information about objects' movement,
position, shape, and contours.  We \textit{do} exercise deliberate
control over such data by conscious choices of where to
direct attention.
} % end paragraph 

\p.1{%
Presumably this balance between conscious and subconscious
cognition is a fine equipoise, evolved to optimize the
efficacy of conscious control.  But even if perceptual
processes such as textural segmentation, occlusion compensation,
and gradient recognition are not guided by conscious
intervention, it is only \textit{through} consciousness
that such gestalts factor in to larger-scale
human actions.  The human mind works in such a manner
that pre-conscious operations become consequential
as they are woven into conscious experience,
against which we orient as perceiver to perceived,
while also superimposing planning and executive
control, situational awareness, personal identity,
and our overall commitment to engage the world
in rational, purposeful ways.  Consciousness
is a necessary medium through which the
\q{outputs} of cognitive processes acquire the status
of perceptual givens, that we deem to disclose
external states of affairs, relative to which we coordinate actions.
} % end paragraph 

\p.1{%
Any comprehensive theory of human cognition, then, is necessarily a
theory of how subjective experience serves as a unifying
field for preconscious faculties, connecting momentary
perceptions to larger-scale goal-directed action and
situational reasoning.  Consequently, investigations
that abstract away an encompassing first-personal context
are only provisional \mdash{} the causal efficacy of a single
nerve cell, neuronal cluster, brain region, or
functional algorithm replicated via AI, depends on
how such entities or processes become manifest in
consciousness.  Biologists notice a similar contrast
between \textit{in vitro} and \textit{in vivo}: molecular
interactions (proteins, etc.) in a test tube versus
in living organisms.  Lab experiments can produce useful
insights, but are simplifying abstractions.  I propose
the analogy: consciousness is the \textit{in vivo} of human
reality, in that any biological or sociological/economic
principles invoked to explain human actions are not explanatorily
complete until we consider their presence in
first-person consciousness.
} % end paragraph 

\p.1{%
The \textit{in vivo} realm raises metascientific
questions: we cannot dissect, nor neuroimage, an animal in
natural settings.  The \textit{in vitro} is artificially
staged to facilitate data acquisition.  Insofar as
consciousness is the humanities' \textit{in vivo}, we similarly
cannot observe another's mental states as objective givens.
This is a well-known philosophical controversy.  Phenomenology
offer techniques to adjudicate claims about cognition
\textit{in consciousness}, compensating for the private nature
of any one person's experience.  In particular, while introspective
consideration of mental states \textit{as they occur} is imperfect
(analysis alters the analyzed), we do have capabilities to
reflect on subjective episodes retroactively.  Our minds are primed to learn
from experience, and to make inferences about others'
mental states.  Usually this would occur
without explicit systematization, perhaps to some
extent subconsciously, but it is not a great leap
to reflect on past experience with a more
analytic/expository intent.
} % end paragraph 

\p.1{%
Suppose I hold a colorful glass paperweight, maybe some
four centimeters in diameter.  I will most likely shift my
visual focus from place to place a couple times, even when holding the
weight still, to discern the full color patterns on its visible
side.  Then I might turn it over.  Perhaps, from touch, I infer that
the bottom side has a flattened base, which I now confirm visually
while rotating the object manually.  As I just begin that action I
have some anticipation of the obverse side's appearance, which (let's
say) is borne out by what I see next.  So a proto-image of that
hidden side is protained while the top is visible, and when the
bottom is visible instead my perception of the top is retained.
The interplay of protention and retention creates a
perceptual synthesis.  Validation of protentive expectations
plays a logical role in such synthesis, as would sense-data
that deviate from anticipated \mdash{} if, say, the paperweight
had a plastic or cloth covering on its base, while I
presumed it were all glass.  Our ability to both formulate
perceptual expectations and adjust to subsequent sense-data
gives a rational structure to the minimal exploratory
patterns of sense-perception.
} % end paragraph 

\p.1{%
Via discussions such as these we can collectively analyze
phenomena such as retention, protention, and perceptual
synthesis.  Empirical actuality
is not required, nor even particularly relevant, for the
analysis.  I am not literally holding a paperweight while
writing these paragraphs, and a paperweight does not
ship with this paper for experimentation.  I assume that
both myself and any reader can engage in such hypothetical
considerations, and two people may either agree or question
each other's use of terms and concepts such as
protention and retention.  This gives the analysis
a public-discursive face, even without empirical
science as normally understood.  We aren't using
EKG machines or other brain-scans to measure neuron activity
while anyone examines a paperweight.  There may or
may not be identifiable neurological correlates to the
concepts of (e.g.) protention and retention.
Searching for such correlates would be a valid investigation
bridging the (metaphorical) \q{in vivo} and \q{in vitro}
facets of human minds.  But as a high-level science,
cognitive phenomenology does not need to develop a
complete theory about how the parameters in its
explanatory arsenal emerge from physical substrata.
It's goal, rather, should be to work toward
consensus about which theoretical ideas and
posits have analytic value and, among such
foregrounded concepts, strive for consistency
in theory-semantics and how analytic parameters are
structurally employed for argumentation.
} % end paragraph 

\p.1{%
Phenomenologists from Husserl onward have built a
literature of analyses and nomenclature toward such
theoretical ends.  Phenomenology explores how minimal sensate
givens are contextualized in situations and subjective
tableaux.  Apperceptive synthesis
occurs on multiple time scales, with any momentary perception
connected to immediate past and future instants via
retention and protention.  We cannot posit
consciousness within single moments in time; the
minimal units of experience, evidently, are brief
but extended perceptual episodes (one or several
seconds' duration, say).  Such episodes then aggregate into
larger units whose inner boundaries, we might propose,
are defined by deliberate acts of perceptual reorientation
(as in the example of switching from watching for an
arriving train to focusing toward its window decal).  And such
compound episodes chain together into purposeful, situational
action, together with the observations and
conceptualizations that sustain it.
} % end paragraph 

\p.1{%
We therefore have a provisional network of theoretical
concepts \mdash{} qualia, proprioception, protention, retention, synthesis,
episodicity, situationality, purposefulness, intentionality
\mdash{} that help organize discussion about conscious states
and attitudes.  Phenomenology
does not restrict itself to sensory introspections
like \q{this shade of red looks identical to that one};
instead, we can reflect on our own experiences to
consider how concepts such as intentionality and
episodicity form building blocks of rigorous
cognitive models.
} % end paragraph 

\p.1{%
Ultimately, we might posit a reductive stack leading from
such subjectively regulated analyses of consciousness down to
brain regions, neuronal complexes, and individual nerve
cells.  Like many inter-theory reductions, there may be
an aura of mystery about the middle stages of this diagram;
perhaps amplified here by the ontological contrast between
conscious states, on the one hand, and physical cells,
on the other.  But the metaphysical controversies
surrounding subjective states should not preclude
a discursive rationality or disputational metatheory for collectively
assessing claims about cognitive-subjective
structuration, even if the warranting data for
any one such claim are not publicly observable.
My argument is that by systematizing a methodology
for analysis of cognition \q{in-consciousness}
\mdash{} by analogy to \q{in vivo} versus \q{in vitro}
\mdash{} phenomenology thereby establishes methodological
prototypes for humanistic and social sciences in
general, at least for that part of their
domains where issues of human thought and decision-making
are relevant.  Phenomenological terminology is
central to such goals; its mode of analysis is not
open-ended, but grounded on the deployment of
specific ideas with relatively fixed technical
meanings.  Via this sense of discursive \q{practice} I
will argue for connections between phenomenology
and cognitive linguistics.
} % end paragraph 


\s|2|{}

\p.1{%
Similar to the controversial status of first-person reports
vis-a-vis theoretical confirmability/falsifiability, the
linguistics community is divided over the value of
analyses based on subjective assessments of acceptability.
It is, of course, common for linguists of many persuasions
to construct ad-hoc sentences/utterances in service
of theoretical claims.  Such methods
raise two questions: whether samples not taken from
real-life corpora are empirically valid, and whether
\mdash{} testing that sentences is meaningful or grammatical
\mdash{} language-user can rely on just their own judgment
(rather than conducting some survey) as conclusive.
} % end paragraph 

\p.1{%
There are interesting analogies between the epistemic
relevance of subjective accounts in linguistics
and in phenomenology.  Both disciplines
have well-rehearsed protocols even while normal
empirical disputation relative to publically-accessible data
does not apply.  The possibility of one writer assenting
or objecting to another's formulation implies that
claims largely endorsed by a discourse community
serves witness to their theoretical value \mdash{} no
less than claims warranted by more empirical measures.
} % end paragraph 

\p.1{%
In one discussion that elucidates metascientific
background for syntactic theory, Hubert Kowelewski
presents the work of Ronald Langacker, and by extension
Cognitive Grammar overall, in terms of \q{Structural Empiricism}:
} % end paragraph 

\p.1{%
Viewed through the lens of structural empiricism, Cognitive Grammar
is therefore a family of models of a linguistic system.  Like any other scientific theory,
it features theoretical entities and processes (cognitive domains,
constructional schemas, compositional paths, subjectification, etc.)
and empirical substructures isomorphic (in a weaker version: similar) to appearances
of observable phenomena. (p. 7)
} % end paragraph 

\p.1{%
He immediately avows \q{It is somewhat unclear what should count as observable phenomena in
linguistics}, but situates such uncertainty in a larger philosophical context.
In particular, \textit{most} scientific theories entertain taxonomies and analyses of
a domain of entities encompassing both public observables and \q{hypothesized}
posits that \mdash{} for one of several possible reasons \mdash{} cannot be \q{seen}.
Empiricism is \textit{structural} insofar as such \q{unobservables} can be
taken seriously because they contribute to persuasive explanations
of \textit{observables'} behavior, with structural models unifying the seen
and the hidden into a rational package.  Sometimes unseen posits are
deemed just as real as everything else, merely beyond our capabilities
of (instrument-enhanced) observation, at least for now.  Other phenomena
(like \q{dark energy}) are left more open-ended, possibly subsumed under
such ontological realism or possibly artifacts of some more complex process
(like \q{phonons}, which behave like \q{particles} of sound, but
that's largely a metaphor).  Natural selection or \q{design} would not be an
object, but rather an encompassing system wherein adaptation drives
evolution.  Potentially \q{dark} matter and energy will ultimately
be given a similarly indirect and holistic cosmological interpretation,
rather than deemed to refer to a specific kind of field or particle.
} % end paragraph 

\p.1{%
The status of mental attitudes and subjectivity vis-a-vis physics and material
science is more circuitous than these other more-physical \q{unobservables},
but Kowaleski implies that we should treat the cognitive as adjunct to the
physical in a similar structural vein: if explanations quantifying
over subjective states provide convincing accounts of observed
behavior (in this case, linguistic performance), then the first-personal
realm should be judged scientifically legitimate as a kind of \q{hidden variables}
dimension.  Specifically in regard to ad-hoc sentence examples, he argues as follows:

\begin{blockQuote}
The linguistic material analyzed in many ... texts [are] constructed
... to illustrate points under analysis.
For this reason, the analyses can be easily attacked on the grounds that
they are somehow unreliable due to the inauthenticity of data. ...
Many linguists seem to embrace the conviction that linguistics should
study only actual, real-life events of language use, which do not belong anywhere on the above
counterfactuality scale [viz., logical impossibility is much stronger
than the contingent fact of something not having yet occurred, with
various intermediates].  However, this conviction overlooks the fact
that many useful scientific models involve clearly counterfactual scenarios
[whose] role is not merely heuristic or pedagogical. \citePage{kow}(99)
\end{blockQuote}

\nip
He then analyzes how \upAdj{9pt}

\begin{blockQuote}
fictional models in natural sciences can be accounted for in
philosophical terms.  In ... structural empiricism, isolated
systems, frictionless pendulums, and perfect liquids are treated just like
any other theoretical objects, except their theoreticity does not follow from
unobservability, as is the case with electrons and electromagnetic fields, but
from their fictional character.  [A] structural empiricist is not
committed to believing in the actual existence of real-world objects corresponding
to theoretical concepts, but only to the empirical adequacy of the theory
featuring these concepts.  \citePage{kow}(111)
\end{blockQuote}

\nip
A reasonable paraphrase vis-a-vis constructed sentences is that, for purposes
of expositing semantic or syntactic claims, the universe of linguistic
objects is broader than merely the totality of sentences humans have in
fact created.  Ad-hoc hypotheticals also belong to that
domain if they belong to model-building regimes.
} % end paragraph 

\p.1{%
I would put it thus: what makes a sentence an object in the
domain of entities spotlighted by linguistics is not the
fact of its being empirically an example of real language-use, but
rather how upon that sentence we can apply observations, analyses,
annotations, etc. whose ordinary ground is real-world
sentences.  We can, for instance, note syntactic, morphological, or lexical patterns,
and point out where pragmatic interpretation is needed.  Constructed
sentences may or may not conform to English grammar,
but at the question of whether they do so is well-posed.
We can write down each word's part of
speech; track pronoun-antecedent pairings and singular or plural alignment; observe
how inflections indicate verb-tense; and etc.: minimal operations in the
syntactic and semantic realms that apply equally
whether samples are real or hypothetical.
} % end paragraph 

\p.1{%
My earlier subway-platform speculation was similarly about hypothetical
circumstances.  I have caught New York subways in real life,
but that's beside the point: we are primed to learn from experience,
but in a manner that abstracts nonsalient details.
My presentation was not free-associative: via posits
like attention, episodicity, and situationality I bore a matrix of
concepts via which to, in effect, \q{annotate} elements of consciousness.
Instead of focusing on the contrast between real and hypothetical
experiences, methodology should consider whether some target of
analysis is amenable to the notions and structures comprising an
analytic protocol.  A hypothetical object, space, or situation becomes
an instance of a theoretical model via the applicability
that model's operative and descriptive maxims.  The proper explanatory
domain of a theory is any structure that can be analyzed
according to the terms of that theory \mdash{} including artifacts
conjured ad-hoc precisely to test, clarify, or refine the theory.
By analogy, again, complex numbers, despite
lacking physical interpretations (such as spatial distance)
are objects of mathematical reasoning because they are
amenable to many quantitative methods (arithmetic
operations, path-integration, infinite summation, and so forth).
} % end paragraph 

\p.1{%
Constructed examples are prominent in both linguistics
and phenomenology.  Relatedly, the \q{demarcation problem}
\mdash{} defending them as legitimate inquiry, not pseudoscience
\mdash{} plays similarly in both contexts.  Suppose I want to examine a
pair of hypothetical examples like:

\begin{exsGroup}

\exsItem{}   John poured water onto the stove.

\exsItem{}   John spilled water onto the stove.

\end{exsGroup}

\nip
The most obvious point would be how \q{poured} versus \q{spilled} frame the
action in different ways (deliberate or accidental).  By selecting
one or the other, the speaker implicitly expresses beliefs about
John's mental state.  So, choice of words can encode
information about the speaker's epistemic starting-point; the
sentence is not merely notating observable details.
} % end paragraph 

\p.1{%
But let's step back.  For sake of discussion, I will assume \exsRef(15) and \exsRef(16)
occur in contexts where listeners know the speaker saw John firsthand
(the evidence is direct observation, not hearsay).  Implicit
in a communicative claim to \textit{describe} something witnessed
is having been in a perceptual orientation to \textit{see} it.  The
sentences then encode salient features of this orientation:
presumably the speaker was standing/sitting in a position
where both John and the stove were in their visual frame,
and they could focus attention on one or the other as
called for.  Likewise, probably they were talking to
John or in some other way clued in to the ambient situation, so
as to rationally infer John's intentions.  Much of the
\textit{configuration} of this cognitive-perceptual orientation,
thereby summarized, reflects my sketched \q{structural model}
for phenomenological description: the intentional relations toward
objects insofar as we deliberately orient toward them
as a focus of attention; the coexistence of multiple
attentional sites (e.g., John and the stove) with focus
shifting back and forth in larger synthesizing episodes;
the role of larger-scale situational awareness guiding
our granular decisions to direct attention hither and thither.
} % end paragraph 

\p.1{%
And yet, such configurations are also implicit in the
sentences' meanings.  Foci of attention tend to become
signified as nouns; so in both cases we have "John" and
"the stove".  It is John who initiates the action that
affects the stove's state, so the former becomes encoded
as subject and the latter as object.  Since speakers
implicitly expresses the evidence for propositions
stated (or implied) in their discourse, understanding a
sentence is partly a matter of inferring mental states,
which includes perceptual comportment toward
its noun-topics.  Anyone hearing \exsRef(15) and \exsRef(16) will,
at least in some schematic manner, reconstruct a
configuration of attentional focus and situational
construal that I tried to explicate \q{phenomenologically}.
Syntactic categories like nouns and verbs are
rooted in listeners' need to infer speaker's cognitive states.
} % end paragraph 

\p.1{%
As a feature of our
reason and lifeworld, mental states exist
in a vivid way in our own momentary consciousness,
but they also exist in an approximative manner
vis-a-vis people (and often animals) around us.
Others' cognitions being, modulo such imprecision,
a natural feature of our own rationality undergirds
\mdash{} so I claim \mdash{} the feasibility of both
phenomenological and linguistic analysis.
We have a
structural model for how people comport themselves
to situations so as to observe and act rationally;
therewith we seamlessly transition
from nonverbal orientation toward surrounding
environments to framings of states of affairs
in communicable ways, then actual speech-production.
Basic organizational principles migrate from one to
the other, such that cognitive-perceptual structures
(attentional foci, say) become manifest,
within communicative intent and staging, as
linguistic phenomena (e.g., that to which
I can perceptually attend becomes a noun).
} % end paragraph 

\p.1{%
If this is correct, then posits of linguistic argumentation
have origins in cognitive faculties that can be
phenomenologically investigated (and should be if we
want an encompassing, \q{in vivo} theory).  At the
same time, language elucidates
how models of other minds are inherent to our
rational agency; they aren't esoteric superstitions
of a phenomenological sect.  Language is the
best case-study for how phenomenology simply
codifies \q{folk theories} of mind and cognition.
} % end paragraph 

\p.1{%
Phenomenology and cognitive linguistics, then, to
my mind reinforce each other; and each can be said
to depend on the other, from an expository
and justificatory point of view.  More to the current
point, hypothetical situations are methodologically
valid because both disciplines curate a
toolkit of terms and concepts whose
applicability transfers over from constructed
to empirical examples.  And, in particular, the
menagerie of theoretical posits in
phenomenology and in linguistics are interrelated,
as \exsRef(15) and \exsRef(16) hopefully illustrate.
} % end paragraph 

\end{document}